\subsection{Changing the allowed level retains the actual ordinary obstruction}

\begin{corollary}[KT-ORDINARY-SYM2-ALL-LEVELS]
Let $\rho_t$, $e$ and $\iota$ be as in
KT-ORDINARY-SYM2-ADM3. For every positive integer $N$,
\[
                  \mathrm{Adm3U\_s}(\iota,e,N,\rho_t)
\]
holds with the same representation. Its trace difference from
$\varepsilon_\ell^e+1+\varepsilon_\ell^{-e}$ is nonunit
at every element of $G_E$, exactly as at $N=2$.
This statement concerns admissibility; it asserts neither
the predicate GD nor membership in every possible existential
comparison set for the changed parameter $N$.
\end{corollary}
\begin{proof}
In the pinned source,
\texttt{BlochKatoSelmerVanishes.lean}, lines 66--67,
\texttt{KillsMu}$(N,\xi)$ means that $\xi$ is trivial on
every automorphism fixing all $N$th roots of unity. It is
not defined by saying that the order of $\xi$ divides $N$.
The constant character $1$ satisfies this exact predicate
for every $N$.

All global finite characters used in the proof of
KT-ORDINARY-SYM2-ADM3 are constant $1$. The same choices
therefore verify \texttt{LevelAway} and \texttt{LevelAt}
for every $N$: their other clauses, namely continuity,
unramified inertia, the actual ordinary flag, the ordered
cyclotomic inertia powers, and nonadjacency, are unchanged.
The source locators are lines 89--92, 102--107 and
111--114 of that file. The extra \texttt{Th.Unl}
predicate, \texttt{TheoremsMuLinked/AutP.lean}, lines
10--13, uses \texttt{FlagAt}; in
\texttt{AutomorphicForms/LocIsotypic.lean}, lines 31--34,
its only occurrence of $N$ is again \texttt{KillsMu}$(N,\xi)$.
Thus the same flag and trivial characters verify it.

The irreducibility and automorphy clauses in
\texttt{DenseThresholdCases/AfStmt.lean}, lines 29--32,
do not contain $N$. The definition of \texttt{Adm3U\_s}
in \texttt{NumberU3US.lean}, lines 27--29, only adds the
unlinked predicate already checked. This proves the stated
admissibility for every positive $N$. Finally the global
residual semisimplification of $\rho_t$ is independent
of $N$, so (OS.4) gives the asserted nonunit trace difference.
\end{proof}

This rules out a specific level-variation shortcut: allowing
a different level while demanding a unit trace difference
for every representation in the entire admissible class
does not remove these systems. Changing the field, imposing
a genuinely smaller comparison class by an additional
source theorem, or controlling extension modules is a
different arithmetic operation. No conclusion about those
operations follows from this corollary.
